% SPDX-License-Identifier: MIT-0
\documentclass[11pt]{article}

\usepackage[T1]{fontenc}
\usepackage[utf8]{inputenc}
\usepackage{lmodern}
\usepackage[a4paper,margin=2.4cm]{geometry}
\usepackage{amsmath,amssymb}
\usepackage{microtype}
\usepackage{booktabs}
\usepackage{longtable}
\usepackage[dvipsnames]{xcolor}
\usepackage{titlesec}
\usepackage[colorlinks=true,linkcolor=MidnightBlue,urlcolor=MidnightBlue,
            citecolor=MidnightBlue]{hyperref}

\hypersetup{%
  pdftitle={A manifest of the research reports in docs/reports},
  pdfauthor={Vladimir Reshetnikov},
  pdfsubject={Catalogue of 64 independent mathematical research packages}}

\definecolor{Ink}{HTML}{16324F}
\definecolor{Accent}{HTML}{4E5C3D}

\titleformat{\section}{\Large\bfseries\color{Ink}}{\thesection}{0.7em}{}
\titleformat{\subsection}{\large\bfseries\color{Accent}}{\thesubsection}{0.7em}{}

\newcounter{entrynum}
\setcounter{entrynum}{0}

\emergencystretch=1.6em

\def\UrlBreaks{\do\/\do\-\do\_\do\.}
\def\UrlBigBreaks{\do\/}

% \entry{title}{directory}{main document}{source archive}{description}
\newcommand{\entry}[5]{%
  \par\medskip\noindent
  \refstepcounter{entrynum}%
  \textbf{\theentrynum.\; #1}\par\nopagebreak
  \nopagebreak{\raggedright\small\color{Ink}%
    \nolinkurl{#2/}\textbf{\nolinkurl{#3}}\par}\nopagebreak
  \nopagebreak{\raggedright\footnotesize\color{gray}%
    from \texttt{#4}\par}\nopagebreak
  \smallskip\noindent #5\par}

\newcommand{\dir}[1]{\texttt{#1}}

\begin{document}

\begin{center}
  {\LARGE\bfseries\color{Ink} A manifest of \texttt{docs/reports}}\\[0.6em]
  {\large Sixty-four independent mathematical research packages,\\
   classified by subject}\\[1.0em]
  {\normalsize Vladimir Reshetnikov}\\[0.3em]
  {\normalsize 19 September 2026}
\end{center}

\vspace{0.8em}

\section{Scope and conventions}

This document catalogues the sixty-four research packages stored under
\dir{docs/reports}.  Each package was delivered as a single ZIP archive; the
archives have been unpacked, one directory per package, and grouped into the
subject hierarchy described in Section~\ref{sec:tree}.  The archives themselves
are no longer kept in the repository.

Every package is self-contained and follows the same shape: a typeset article
(\texttt{*.pdf}) with its complete \LaTeX{} source, a \texttt{README.md}
summarising the result, usually a \texttt{code/} directory with exact-arithmetic
verification scripts, a \texttt{data/} or \texttt{results/} directory holding the
recorded output of those scripts, and often \texttt{STATUS.md},
\texttt{sources.md} or \texttt{notes/} files auditing the sources consulted and
the precise limits of the claim.

Two cautions apply uniformly and are therefore not repeated in the individual
entries.  First, these are AI-assisted research drafts: none has been
independently refereed, and none has been checked in a proof assistant.  Second,
the descriptions below record \emph{what each report sets out to settle and what
it claims to prove}.  Compiling this manifest involved reading each package's
statement of its problem and result; it did not involve checking any proof.
Where a report itself flags a limitation --- a prior disproof, an imported
theorem, a correction to the source, or an exception left open --- that is noted,
because it changes what the report is.

The naming of the directories is editorial.  Several distinct reports attack the
same published problem, and some arrived under colliding archive names, so each
directory has been given a name describing its own angle on the problem.  The
original archive name is recorded with every entry.

\clearpage
\section{The directory tree}
\label{sec:tree}

\begin{center}
\begin{tabular}{@{}llr@{}}
\toprule
\textbf{Category} & \textbf{Subcategory} & \textbf{Reports}\\
\midrule
\dir{ordinals-and-order-types}    & \dir{friendly-order-types}          & 4\\
                                  & \dir{wqo-powersets-and-statures}    & 4\\
                                  & \dir{transfinite-words}             & 3\\
                                  & \dir{ordinal-arithmetic}            & 2\\
                                  & \dir{games-on-ordinals}             & 2\\
\midrule
\dir{hankel-determinants}         & \dir{somos-and-elliptic}            & 3\\
                                  & \dir{catalan-and-ballot}            & 2\\
                                  & \dir{arithmetic-hankel}             & 2\\
                                  & \dir{growth-and-runs}               & 2\\
\midrule
\dir{congruences-and-valuations}  & \dir{supercongruences}              & 3\\
                                  & \dir{iterated-series}               & 2\\
                                  & (top level)                         & 3\\
\midrule
\dir{tetration-and-digit-stabilization} & ---                           & 7\\
\midrule
\dir{log-concavity-and-unimodality}     & ---                           & 7\\
\midrule
\dir{graph-theory}                      & ---                           & 3\\
\midrule
\dir{enumerative-combinatorics}         & ---                           & 8\\
\midrule
\dir{generating-functions-and-asymptotics}
                                  & \dir{binary-product-coefficient-counts} & 4\\
                                  & (top level)                         & 3\\
\midrule
\multicolumn{2}{l}{\textbf{Total}} & \textbf{64}\\
\bottomrule
\end{tabular}
\end{center}

\medskip
\noindent
More than a quarter of the collection consists of counterexamples and
corrections rather than proofs.  Entries 14, 15, 29, 32, 33, 38, 39, 40, 41, 43,
48, 52, 58, 59, 60, 62 and 64 refute, or replace with a proved statement, an
assertion published or catalogued elsewhere; entry 34 disproves one of the three
conjectures it treats while proving the other two.  Several reports whose main
outcome is a proof still correct the statement they prove: entries 17, 19 and 24
record a sign error, a missing one-unit shift and a numbering ambiguity in their
sources, and entry 44 refines its target after finding that the original
conjecture had already been disproved by someone else.

\section{Ordinals and order types}

\subsection{Friendly order types}

Four independent reports on Isa Vialard's \emph{friendly order type} $f(P)$, the
rank of the tree of open-ended bad sequences in which every chosen point has an
incomparable partner in the current residual.  All four prove the same finite
formula; they differ in what they build on top of it.

\entry{Friendly Order Type and Incomparability Components}
{ordinals-and-order-types/friendly-order-types/incomparability-components-and-forests}
{article.pdf}{friendly\_order\_type\_research.zip}
{Proves $f(P)=|P|-c(\mathrm{Inc}(P))$ for a finite poset $P$, where $c$ counts
connected components of the incomparability graph (isolated vertices included).
With Vialard's width-transfer theorem this gives ordinal width
$\omega^{f(P)}$ for the Dershowitz--Manna multiset ordering.  Also classifies
all maximum friendly subsets, builds spanning-forest certificates, and gives
Cartesian-product, lexicographic-substitution and well-ordered-sum formulas
with an enumerative generating function.}

\entry{Friendly Order Types Through Incomparability Components}
{ordinals-and-order-types/friendly-order-types/incomparability-components-and-infinite-tails}
{friendly\_order\_types.pdf}{friendly\_order\_type\_research (1).zip}
{The same component formula, together with an exact friendly rank for connected
infinite well-partial-orders: writing $o(C)=\lambda+n$ with $\lambda$ the limit
part and $T(C)$ the finite set of points with finite principal upsets.
Includes an exhaustive finite verification and a proof audit.}

\entry{Friendly Order Type Through Incomparability Components}
{ordinals-and-order-types/friendly-order-types/matroid-rank-and-ordinal-sums}
{article.pdf}{friendly\_order\_type\_research (3).zip}
{Identifies $|P|-c(P)$ with the graphic-matroid rank and the incidence-matrix
rank of the incomparability graph.  Adds a substitution formula, a finite
Cartesian-product formula, and additivity over arbitrary well-ordered ordinal
sums; evaluates every well-partial-order whose incomparability components are
finite, and gives countable counterexamples delimiting stronger infinite
extrapolations.}

\entry{Friendly Order Types of Finite Posets and Ordinal Products}
{ordinals-and-order-types/friendly-order-types/cartesian-products-of-ordinals}
{friendly\_order\_types.pdf}{friendly\_order\_types\_research.zip}
{Beyond the component formula: a Cartesian product with at least two
non-singleton factors has exactly one non-endpoint incomparability component;
the tuple $(o,h,w,f)$ of the factors does \emph{not} determine $f(A\times B)$
(explicit four-point $D,E$ with equal tuples and $f(D\times 2)=5$,
$f(E\times 2)=6$); and a complete classification of
$f(\alpha_1\times\cdots\times\alpha_r\times B)$ as a natural product, with an
exception when every $\alpha_i$ is a successor and $B$ has a greatest element.}

\subsection{Well-quasi-ordered powersets and statures}

Four reports on the open questions in Section~6 of Abriola, Halfon, Lopez,
Schmitz, Schnoebelen and Vialard, \emph{Measuring well-quasi-ordered finitary
powersets} (arXiv:2312.14587v2), and on the corresponding topological invariant.

\entry{Cofinal strata and ordinal absorption}
{ordinals-and-order-types/wqo-powersets-and-statures/cofinal-strata-of-finitary-powersets}
{article.pdf}{cofinal\_strata\_ordinal\_powersets.zip}
{For $S_\alpha(Q)$, a finite poset $Q$ with each point replaced by a chain
$\alpha=\omega^\rho$, proves
$o(K(S_\alpha(Q)))=\sum_k \omega^{\rho\otimes k}c_k(Q)$ and
$h(K(S_\alpha(Q)))=\alpha\cdot\mathrm{height}(M(Q))$, where $c_k$ counts size-$k$
maximal antichains not dominated by a larger one.  The maximal-order-type
theorem allows nonuniform pure well-partial-order fibers and is proved by a
finite cofinal-stratification argument.}

\entry{Intrinsic ordinal approximations for finitary powersets}
{ordinals-and-order-types/wqo-powersets-and-statures/intrinsic-ordinal-approximations}
{article.pdf}{intrinsic\_ordinal\_approximations.zip}
{Studies the intrinsic invariant
$a_o(A)=\sup\{o(B)+1: \mathrm{Idl}(B)\hookrightarrow A\}$.  Proves the universal
bound $a_o(A)\le h(P_f(A))$ with equality exactly when every finitely generated
downset of $A$ is an algebraic dcpo, gives the closed formula
$h(P_f(\prod_i\alpha_i))=\sup(\bigotimes_i(1+a_i)+1)$, and shows the universal
identity fails on $\lambda+F$.}

\entry{Ordinal Order Maps and Finitary Powersets of Ordinal--Finite Grids}
{ordinals-and-order-types/wqo-powersets-and-statures/ordinal-order-maps-and-grids}
{ordinal\_order\_maps.pdf}{ordinal\_order\_maps\_research.zip}
{Exact maximal order types and heights for finite Hoare powersets of
$\beta\times Q$ with $Q$ a finite poset: for
$\beta=\omega^{\gamma_0}+\cdots+\omega^{\gamma_{m-1}}$,
$o(\mathrm{Mon}(Q,\beta))=\bigoplus_{s\,\text{isotone}}\omega^{\bigoplus_q\gamma_{s(q)}}$.}

\entry{Sharp ordinal bounds for Hoare powerspaces}
{ordinals-and-order-types/wqo-powersets-and-statures/hoare-powerspace-statures}
{article.pdf}{hoare\_statures\_research.zip}
{For all Noetherian spaces, $1+\|X\|\le\|HX\|\le 2^{\|X\|}$, with witnesses
attaining both endpoints at every ordinal and $T_1$ upper witnesses; the older
$\omega^\alpha$ bound is attainable at infinite input stature $\alpha$ exactly
when $\omega\cdot\alpha=\alpha$; the complete spectrum at $\omega+n$ is
$\omega\cdot j(U)+(j(Q)-j(U))$.  Extends the known well-quasi-order results of
Abriola et al.\ to arbitrary Noetherian spaces.}

\subsection{Transfinite words}

Three reports on Harry Altman, \emph{Bounding finite-image sequences of length
$\omega^k$} (arXiv:2409.03199v2), Example~3.15: is the bound
$\omega^{\omega^4}$ sharp for words of length ${<}\,\omega^2$ over the
two-element antichain?

\entry{Exact maximal order types of finite-alphabet transfinite words}
{ordinals-and-order-types/transfinite-words/finite-alphabet-transfinite-words}
{article.pdf}{finite\_alphabet\_transfinite\_words.zip}
{Proves $o(W_k(P))=\omega^{\omega^{N_k(P)-1}}$ for words of ordinal length
${<}\,\omega^k$ over a finite poset $P$ under subsequence embedding, where
$N_k$ counts indecomposable classes of a canonical finite atom poset
$I_k(P)$, with $N_1(P)=|P|$ and $N_{k+1}(P)=|P|+J(I_k(P))-1$.  For $k=2$ the
binary antichain gives $\omega^{\omega^4}$ --- the tightness Altman leaves
open --- and the binary chain gives $\omega^{\omega^3}$.}

\entry{Exact maximal order types below omega squared}
{ordinals-and-order-types/transfinite-words/order-types-below-omega-squared}
{article.pdf}{omega2\_order\_types\_proposed\_solution.zip}
{The $k=2$ case in its own right:
$o(s^F_{\omega^2}(P))=\omega^{\omega^{n+j(P)-2}}$ for $n=|P|\ge 2$, with $j(P)$
the number of downsets including the empty one.  Handles the empty and
singleton alphabets as genuine exceptions and classifies arbitrary selected
families of periodic tail types.}

\entry{Guarded periodic blocks and maximal order types below omega squared}
{ordinals-and-order-types/transfinite-words/guarded-periodic-blocks}
{article.pdf}{ordinal\_words\_research.zip}
{The same formula $o(s_{\omega^2}(X))=\omega^{\omega^{n+|J(X)|-2}}$, obtained by
a guarded-periodic-block construction, with symbolic algorithms and executed
tests; the one-letter alphabet (type $\omega^2$) is flagged as outside the
formula's range.}

\subsection{Ordinal arithmetic}

Two independent proposed solutions to Problem~6.2 of Paolo Lipparini,
\emph{Monotone infinitary operations on ordinals} (arXiv:2505.00424v2): define
and study the least weakly monotone, $e$-special strictly increasing operation
on countable sequences of ordinals.

\entry{An Explicit Formula for the Minimal $e$-Monotone Infinitary Ordinal Sum}
{ordinals-and-order-types/ordinal-arithmetic/lipparini-minimal-infinitary-sum}
{article/explicit\_ordinal\_sum.pdf}{lipparini\_problem\_6\_2\_proposed\_solution.zip}
{An explicit formula for the least admissible operation, proved by an arithmetic
admissibility upper bound and a rank-theoretic minimality lower bound, using an
explicit countable-multiset profile order, its rank, absorption and correction
rules, bounded-poset heights and finite-cap ranks.  Imports Lipparini's
formula and minimality theorem for the weaker operation $S$.}

\entry{An Explicit Formula for Lipparini's Minimal Infinitary Ordinal Operation}
{ordinals-and-order-types/ordinal-arithmetic/lipparini-minimal-operation-formula}
{article.pdf}{ordinal\_operation\_research.zip}
{A three-case finite ordinal-arithmetic formula with proofs of admissibility and
pointwise minimality.  Relative to Lipparini's already evaluated $S$, the
correction is zero, a finite increment, or one \emph{natural} summand $\omega$
--- ordinary addition of $\omega$ is not a substitute.}

\subsection{Games on ordinals}

\entry{A counterexample to winner stability in ordinal Chomp}
{ordinals-and-order-types/games-on-ordinals/ordinal-chomp-winner-stability}
{ordinal\_chomp\_counterexample.pdf}{ordinal\_chomp\_counterexample.zip}
{For the numerical semigroup $S=\langle 4,6,9\rangle$, Chomp on $S$ is a
second-player win while Chomp on the ordinal square
$S^2=\{\omega a+b: a,b\in S\}$ is a first-player win --- the opening
$\omega\cdot 4+4$ leaves Grundy value zero, and the same move wins at every
exponent ${\ge}\,2$, so $\mathrm{ch}(4,6,9)=2$.  This refutes Conjecture~4.1 of
arXiv:2504.07317v1 and answers Question~4.6 of v3 negatively, already for
natural-number generators.}

\entry{Open-query games on ordinal spaces}
{ordinals-and-order-types/games-on-ordinals/open-query-membership-games}
{article.pdf}{ordinal\_membership\_games.zip}
{A negative answer to Problem~1.3 of Chiozini, Csern\'ak and Soukup
(arXiv:2510.05754v3): two copies of $[0,\omega]$ each have set-membership
number~1, but their topological sum has value~2.  The general theorem computes
the invariant for every nonempty Hausdorff space of finite Cantor--Bendixson
height~$h$ as $\lceil\log_2(h+1)\rceil$ when the last nonempty derivative is a
singleton and $\lceil\log_2(h+2)\rceil$ otherwise.}

\section{Hankel determinants and continued fractions}

\subsection{Somos sequences and elliptic curves}

\entry{A proof of Barry's Hankel--Somos-6 conjecture}
{hankel-determinants/somos-and-elliptic/barry-somos-6-hankel-recurrence}
{article.pdf}{barry\_somos6\_research.zip}
{Conjecture~7 of Paul Barry, arXiv:2211.12637: for the Hankel determinants
$H_N$ of a four-parameter Jacobi-type generating function,
$H_NH_{N-6}=\alpha H_{N-1}H_{N-5}+\gamma H_{N-3}^2$ with
$\alpha=t^2(r+2)^2$ and
$\gamma=t^3((s+t)(r+s+t+2)^2+t(r+2)^3)$.  Also proves a homogeneous
four-parameter form and a seven-parameter polynomial Hankel theorem making
exceptional specialisations rigorous; the verifier never divides by a
determinant.}

\entry{Finite Hankel Certificates for Three Somos-8 Conjectures}
{hankel-determinants/somos-and-elliptic/somos-8-hankel-certificates}
{article.pdf}{somos\_hankel\_proofs.zip}
{Conjectures~11, 12 and~13 of the same Barry paper, by importing Hone's
genus-two Hankel theorem (arXiv:1907.05204v3, Thm.~5.4) and supplying a
binomial translation, a polynomial-minor specialisation valid on the degenerate
locus, and a rank-bound lemma that turns four symbolic rows into an all-index
identity.  Identifies a sign error in the source's Catalan expression for
Conjecture~13.}

\entry{Elliptic addition and Jacobi continued fractions}
{hankel-determinants/somos-and-elliptic/elliptic-addition-and-jacobi-fractions}
{article.pdf}{elliptic\_continued\_fraction\_research.zip}
{A consistently indexed, non-torsion version of Barry's 2023 Conjecture~4, with
the coordinate formulas of Conjecture~3: every continued-fraction tail is given
explicitly and removing one layer is shown to be elliptic translation.
Hypotheses: characteristic zero, $Y^2+aXY+bY=X^3+cX^2+dX$ nonsingular,
$b\ne 0$, and $P=(0,0)$ of infinite order; order-three and order-four examples
show the regular fraction genuinely needs the last one.}

\subsection{Catalan and ballot moments}

\entry{Shifted Catalan Hankel Polynomials}
{hankel-determinants/catalan-and-ballot/shifted-catalan-hankel-polynomials}
{article.pdf}{catalan\_hankel\_conjecture\_solution.zip}
{For $D_N^{(m)}(a,b)=\det(aC_{i+j+m}+bC_{i+j+m+1})$, proves the
denominator-exponent assertion on page~4 of Barry, arXiv:2011.10827v1, and
shows the separately printed Conjecture~2 needs an explicit one-unit shift
correction.  The minimal characteristic polynomial is
$(X^2-(a+2b)X+b^2)^{m(m-1)/2+1}$, so the minimal order is $m(m-1)+2$; all
exceptional reduced denominators, the exact coefficient formula, numerator
degree and reciprocal symmetry are supplied.}

\entry{Product formulas for ballot-polynomial Hankel determinants}
{hankel-determinants/catalan-and-ballot/ballot-polynomial-hankel-determinants}
{article.pdf}{Cigler\_Hankel\_Conjectures\_13\_15.zip}
{Conjectures 13, 14 and 15 in Section~5 of Johann Cigler, arXiv:2111.14492v3.
An explicit coefficient product gives the sign and exact support of the
normalised shift-$k$ Hankel determinant of the ballot moments and positivity of
every residual coefficient for $n\ge1$; the same products give the conjectured
generating-function denominators, exact numerator degrees and reciprocal
numerator identities.  Adds strict log-concavity of the residual coefficients
and their limiting binomial profile.}

\subsection{Arithmetic Hankel determinants}

Two independent reports on the numerator formula recorded as a conjecture in
OEIS A122251, with A122249 as its binary special case.

\entry{A Proof of an OEIS Hankel-Numerator Conjecture}
{hankel-determinants/arithmetic-hankel/a122251-numerator-formula}
{report.pdf}{hankel\_numerator\_research.zip}
{For $H_n(m,a)=\det(1/(a+m(i+j)))_{0\le i,j\le n}$, the reduced numerator is
$\prod_{p\mid m}p^{\,n(n+1)v_p(m)+2\sum_{k\le n}v_p(k!)}$: independent of the
coprime offset $a$, a perfect square, and supported on the primes dividing $m$.
Also gives exact denominator valuations, the gcd of denominators over all
coprime $a$, a constructive congruence witness, the noncoprime-parameter
formula, recurrences and a natural-boundary proof.}

\entry{Prime-supported numerators of arithmetic Hankel determinants}
{hankel-determinants/arithmetic-hankel/a122251-prime-exclusion-proof}
{report.pdf}{hankel\_oeis\_conjecture\_proof.zip}
{The same formula in the size-$m$ indexing,
$\prod_{\ell\mid d}\ell^{\,v_\ell(d)m(m-1)+2\sum_{j<m}v_\ell(j!)}$, with the
exclusion of every prime not dividing $d$ carried out as a separate step ---
the point being that determining the $p$-adic valuation alone would not
suffice.  Every reduced denominator valuation is determined as well.}

\subsection{Growth laws and run statistics}

\entry{A proof of the Ap\'ery Hankel growth conjecture}
{hankel-determinants/growth-and-runs/apery-hankel-determinant-growth}
{apery\_hankel\_growth.pdf}{apery\_hankel\_research.zip}
{For $A_k=\sum_j\binom{k}{j}^2\binom{k+j}{j}^2$ and $D_n=\det(A_{i+j})$, proves
$\log D_n=n(n+1)\log\frac{17+12\sqrt2}{4}+O(n)$ --- the limit conjectured in the
September 2026 revision of OEIS A228143 --- using Edgar's moment-density
theorem.  Also an explicit law for $\det(A_{m(i+j)+r_n})$ with $r_n/n\to\rho$,
and non-D-finiteness of the formal generating series of $D_n$.}

\entry{Binary Runs in Rueppel Hankel Determinants}
{hankel-determinants/growth-and-runs/rueppel-binary-run-determinants}
{rueppel\_hankel.pdf}{rueppel\_hankel\_research.zip}
{For the Rueppel series $r(x)=\sum_{j\ge0}x^{2^j-1}$, the polynomial identity
$H_n(tx+1/r(x^2))=(-1)^{1+\varepsilon+\lfloor R(n)/2\rfloor}
(\varepsilon+(R(n)-\varepsilon)t^2)$, where $R(n)$ counts runs of equal bits in
the binary expansion of $n$ and $\varepsilon$ is the low bit of
the Gray code $n\oplus\lfloor n/2\rfloor$.  Settles statements 8, 10, 13 and 23 of Barry's 2021 paper
(not the identically numbered ones elsewhere), with a correction.}

\section{Congruences and $p$-adic valuations}

\subsection{Supercongruences}

\entry{A cubic supercongruence for OEIS A348410}
{congruences-and-valuations/supercongruences/a348410-cubic-supercongruence}
{article.pdf}{A348410\_Supercongruence\_Research.zip}
{For $A_{u,v}(n)=[x^n](1-x)^{-un}(1+x)^{-vn}$ proves
$v_p(A_{u,v}(N)-A_{u,v}(N/p))\ge 3v_p(N)-\varepsilon_p$ with
$\varepsilon_3=1$ and $\varepsilon_p=0$ for $p\ge5$, for every odd prime $p$ and
every $N$ divisible by $p$ --- no coprimality condition on the remaining
multiplier.  Taking $(u,v)=(2,1)$ gives Peter Bala's A348410 supercongruence.
Appendix~A audits arXiv:2104.10754v1, whose prime-by-prime framing conclusion
and weighted harmonic assertion fail on its own displayed examples.}

\entry{Small-prime congruences for A348410}
{congruences-and-valuations/supercongruences/a348410-small-prime-congruences}
{article.pdf}{a348410\_supercongruence.zip}
{For $a(n)=[x^n]((1-x)(1-x^2))^{-n}$ proves
$v_p(a(mp^r)-a(mp^{r-1}))\ge 3r-v_p(12)$ for every prime and all $m,r\ge1$,
hence $(12/n^3)\sum_{d\mid n}\mu(d)a(n/d)\in\mathbb{Z}_{\ge0}$ with $12$
optimal.  The principal extension is an exact binary leading-defect formula
for the family $A_{\alpha,\beta}$, valid modulo $2^{3r-2}$.}

\entry{Signed binomial supercongruences and OEIS A352373}
{congruences-and-valuations/supercongruences/a352373-signed-binomial-supercongruences}
{signed\_binomial\_supercongruence.pdf}{OEIS\_A352373\_supercongruence.zip}
{For $a(n)=[x^{tn}](1+x)^{rn}(1-x)^{sn}$ proves
$a(mp^h)\equiv a(mp^{h-1})\bmod p^{3h}$ for $p\ge5$ and
$a(m3^h)\equiv a(m3^{h-1})\bmod 3^{3h-1}$, with $m$ not required coprime to
$p$.  Covers A352373, A348410, A351856 and A351857 as specialisations; the only
non-elementary import is the classical Jacobsthal binomial congruence.}

\subsection{Congruences for iterated and compositional series}

\entry{Arithmetic of a Self-Iterating Generating Function}
{congruences-and-valuations/iterated-series/a396807-self-iterating-series}
{article.pdf}{A396807\_congruences\_and\_prime\_power\_lifts.zip}
{Both conjectures in OEIS A396807 for $A(x)=x+A^{[5]}(x)A^{[6]}(x)$ (functional
iteration): $A(x)\equiv x/(1-x)$ and $A^{[k]}(x)\equiv x/(1-kx)$ modulo $10$ for
every integer $k$, negative iterates included.  Adds a sharp modulus theorem
for $F=x+cF^{[r]}F^{[s]}$, rational lifts at every selected prime power, and
minimal period $46860$ for A396807 modulo $100$ starting at index~1.}

\entry{Prime-power congruences for compositional tree series}
{congruences-and-valuations/iterated-series/compositional-tree-series-congruences}
{article.pdf}{compositional\_tree\_congruences.zip}
{For $A_l(x)=x\exp(A_l^{\circ l})$, shows the OEIS A396805 claim
$a_5(n)\equiv n\ (\mathrm{mod}\ 5)$ is \emph{false}: first failure $n=23$, with
exception set $\{20j+3,\,20j+4: j\ge1\}$ and residue~2 throughout.  Proves
eleven other conjectural assertions in A396803--A396806 via a general reduction
of coefficients mod $p^\alpha$ to typed rooted trees on fewer than $p^\alpha$
labels, yielding eventual periods, rational residue generating functions, and
an exact classification of $a_l(n)\equiv n\pmod p$.}

\subsection{Valuations, divisibility and periodicity}

\entry{Binary Carries and the Divisibility of Two-Stack-Sortable Permutation Numbers}
{congruences-and-valuations/a000139-binary-carries-parity}
{research.pdf}{a000139\_carries\_research.zip}
{Proves Peter Bala's parity conjecture for OEIS A000139,
$a(n)=2(3n)!/((n+1)!(2n+1)!)$: $a(n)$ is odd exactly when $n$ is odd and its
binary expansion contains no \texttt{11}.  The full valuation is
$v_2(a(n))=s(n)+s(n+1)-s(3n)$ in binary digit sums.  Two independent parity
proofs, plus exact valuation distributions, recurrences, asymptotics, extremal
formulas and limit laws.}

\entry{Polynomial divisibility of height-one factorial ratios}
{congruences-and-valuations/factorial-ratio-polynomial-divisibility}
{article.pdf}{bala\_factorial\_divisibility\_research.zip}
{A general rational-root criterion for uniform polynomial divisibility of
balanced integral factorial ratios, applied to Peter Bala's uniform-product
conjectures for OEIS A211417 and A295431: explicit multipliers for every product
length, with the least multipliers certified in eight particular examples.
Includes an atlas of the 52 sporadic height-one ratios with recurrences,
generating functions and asymptotic expansions.}

\entry{When the Secant Numbers Become Periodic}
{congruences-and-valuations/secant-number-periodicity}
{article.pdf}{secant\_periodicity\_research.zip}
{Refutes the OEIS A000364 conjecture that the secant numbers from $n=1$ are
purely periodic modulo every $m$: at $m=27$, $a(1)\equiv1$ but
$a(19)\equiv10$.  The exact replacement is that pure periodicity from index one
holds if and only if no odd prime cube divides $m$; the least eventual period
always divides $\varphi(m)$.  Determines the least eventual period and exact
onset for every modulus, with densities and a certified mean onset.}

\section{Tetration and digit stabilization}

\entry{A Counterexample to a Decimal-Length Bound for Tetration Stabilization}
{tetration-and-digit-stabilization/congruence-speed-decimal-length-bound}
{article.pdf}{tetration\_counterexample.zip}
{An explicit 2544-digit base $A=1+qu$ ($q=5^{2547}$) for which
$\nu_{10}(T_{n+1}-T_n)=\min(1+2548n,\,2547(n+1))$ at every height, so the
congruence speed at height $\mathrm{len}(A)+2$ is $2548$ while the permanent
value is $2547$.  This contradicts the decimal-length form of Conjecture~1 of
Marco Rip\`a, \emph{The congruence speed formula} (NNTDM 27(4), 2021).  The
report audits the later Rip\`a--Onnis confirmation: their weaker valuation
bound is compatible with the counterexample, the decimal-length bound is not.}

\entry{Exactly One Digit per Height}
{tetration-and-digit-stabilization/one-stable-digit-per-height}
{article.pdf}{tetration\_exact\_digits.zip}
{The three conjectural comments in OEIS A324017: the first two are proved and
the third is disproved for every nontrivial base in its stated domain.  For
even $q\ge4$ and $a=q-1$, the central theorem is
$T_{h+r}-T_h\equiv-2q^h\pmod{q^{h+1}}$, so each added height contributes
exactly one new stable radix-$q$ digit.  Smallest counterexample to the third:
$3\uparrow\uparrow2\equiv27$ and $3\uparrow\uparrow3\equiv59 \pmod{64}$.
Adds exact prime-adic distances, a factorization-free digit-lifting algorithm
and higher Knuth arrow ranks.}

\entry{Tetration, perfect powers, and frozen decimal digits}
{tetration-and-digit-stabilization/perfect-power-minima-for-frozen-digits}
{article.pdf}{tetration\_perfect\_power\_minima.zip}
{Determines $\rho(n)$, the smallest $n$th perfect power $x^n>1$ whose eventual
number of newly frozen terminal decimal digits is exactly $n$:
$\rho(1)=2$, $\rho(2)=49$, and $\rho(n)=B(s)^n$ for $n\ge3$ with
$s=n-v_2(n)$ and $B(s)\in\{2^s\pm1,\,3\cdot2^s\pm1\}$ according to
$s\bmod4$.  All other roots are ruled out by a deterministic quadratic size
barrier plus small exact certificates.}

\entry{Phase-shift classification for integer tetration}
{tetration-and-digit-stabilization/phase-shift-classification}
{tetration\_phase\_classification.pdf}{tetration\_phase\_classification.zip}
{Proves the 21-word phase-shift conjecture in OEIS A376842, the 14-word
even-base restriction from the appendix of Rip\`a's 2025 paper, and the
28-word prescribed-anchor conjecture in OEIS A376446.  The mechanism compares
the 2-adic and 5-adic valuations of $D_b=T_{b+1}-T_b$ at the first
permanently constant-speed height: strict 2-adic dominance gives phase 5,
strict 5-adic dominance a multiplicative orbit in $\mathbb{F}_5^\times$, and
balanced valuations exactly four further anchored words.}

\entry{A Sharp Convergence Theorem for the $q$-Newton Series of Subcritical Tetration}
{tetration-and-digit-stabilization/q-newton-series-convergence}
{article/tetration\_qnewton.pdf}{tetration\_qnewton\_research.zip}
{Proves convergence of the explicit $q$-binomial interpolation series of
MathOverflow question 259278 for every real base $1<a<e^{1/e}$, with
$\log a=qe^{-q}$: the grouped series converges absolutely and locally uniformly
for $\mathrm{Re}\,z>-2$, converges conditionally on $\mathrm{Re}\,z=-2$ except
at $z=-2+2\pi ik/(-\log q)$, and diverges for $\mathrm{Re}\,z<-2$.  Adds strict
complete monotonicity and uniqueness under an eventual log-convexity
condition.}

\entry{Delayed stabilization of power towers}
{tetration-and-digit-stabilization/delayed-stabilization-of-towers}
{article.pdf}{tetration\_stabilization\_research.zip}
{Refutes the claim that every prime base ending in 9 reaches its permanent
decimal congruence speed by height~2 (Rip\`a, Conjecture~2, journal p.~57;
arXiv:2208.02622v1, Conjecture~4.1): $2749$ is the smallest prime
counterexample, with first permanent-speed height $H(2749)=3$.  Proves exact
formulas for $C_a(h)$, $V(a)$ and $H(a)$ in terms of
$s=v_2(a-1)$, $t=v_2(a+1)$, $u=s+t-1$, $e=v_5(a+1)$, and shows there are
infinitely many primes with each prescribed onset $H\ge3$.}

\entry{Delayed digit stabilization: a counterexample and sharp bounds}
{tetration-and-digit-stabilization/iterated-power-digit-stabilization}
{delayed\_digit\_stabilization.pdf}{digit\_stabilization\_research.zip}
{For $S_a(b)=v_{10}(a^{10^{b+1}}-a^{10^b})$ and $D_a(b)=S_a(b)-S_a(b-1)$,
answers negatively the question whether $D_a(b)=1$ for all $b\ge3$: the
smallest counterexample is $a=3\cdot2^{99}+1$, with $S_a(2)=100$,
$S_a(3)=102$, $D_a(3)=2$.  Proves minimality, an exact onset formula, the
smallest counterexample at every stage, the natural density of exceptions, and
arbitrary-radix extensions.  The article is explicit that $a^{10^b}$ is not
genuine tetration and derives the tower formulas separately.}

\section{Log-concavity, unimodality and positivity}

\entry{Single curves do not force log-concavity}
{log-concavity-and-unimodality/cluster-variable-log-concavity-counterexample}
{article.pdf}{annular\_single\_lamination\_counterexample.zip}
{An unpunctured annulus with one marked point on each boundary and a single
unit-weight simple lamination curve gives, after mutations $2,1,2$ with the two
mutable variables specialised to~1, the polynomial
$1+3q^2+2q^3+4q^4+2q^5+q^6$; the consecutive coefficients $3,2,4$ form a strict
valley, refuting Conjecture~1 of arXiv:2508.04396v3 and conflicting with its
general unpunctured-surface unimodality statement.  The geometric realisation
is proved, not merely the matrix calculation, and a fixed seed gives infinitely
many such examples.}

\entry{Cycles refute infinite log-concavity for chromatic coefficients}
{log-concavity-and-unimodality/infinite-log-concavity-of-cycles}
{article.pdf}{chromatic\_cycle\_logconcavity.zip}
{For the cycle $C_n$ with absolute chromatic coefficients under the iterated
operator $L(a)_k=a_k^2-a_{k-1}a_{k+1}$: infinitely log-concave for
$3\le n\le11$; first failure at depth 5 for $n=12$, depth 4 for
$13\le n\le16$, and depth 3 for $n\ge17$.  So $C_{12}$ is the smallest cycle
counterexample (no minimum over all graphs is claimed), with an explicit
third-iterate witness polynomial valid for all $n\ge3$.}

\entry{Oscillating Log-Concavity in Geometric Generating Functions}
{log-concavity-and-unimodality/oscillating-log-concavity-geometric-transform}
{article.pdf}{oscillating\_log\_concavity\_research.zip}
{Challenges 2 and 3 in Section~3 of Heim and Neuhauser, arXiv:2302.13327v1, for
the \emph{geometric} transform $1/(1-uG(z))$.  For divisor-sum weights, every
real exponent and every $u>0$ give infinitely many strict log-concavity and
log-convexity indices with eventually bounded gaps; likewise for integer power
weights with exponent ${\ge}\,2$, the exponents $0$ and $1$ being explicitly
exceptional.  For square weights at $u=1$, every four consecutive indices carry
both signs (sharp), each with density $1/2$ in every arithmetic progression.
A subdominant-pole theorem is the common mechanism.}

\entry{A counterexample and a sharp spacing-three criterion for products of $q$-integers}
{log-concavity-and-unimodality/q-integer-product-unimodality}
{report.pdf}{q\_unimodality\_attack.zip}
{The necessity half of Conjecture~5.4 of Connelly, Ito, Martinez, Shevchenko and
Yang (arXiv:2605.12822v1) is false: $(1+q)^6(1+q^3)$ is symmetric and weakly
unimodal although no $a_i$ is divisible by 3 and $b=2$ exceeds the asserted
bound, and the example lies inside the source's claimed range.  The report also
proves the proposed sufficient condition for every $r\ge2$, a necessary degree
bound in the absence of an $r$-divisible ordinary factor, and a complete
classification of spacings two and three.}

\entry{Sharp size thresholds for failures of log-concavity in independence systems}
{log-concavity-and-unimodality/independence-system-thresholds}
{article.pdf}{sharp\_log\_concavity\_thresholds.zip}
{For each rank $r\ge4$, determines the smallest ambient ground-set size
admitting failure of ordinary, ordered and ultra log-concavity, and shows the
three thresholds coincide across all independence systems, those satisfying
Xie--Xu's augmentation bound, and those with additive augmentation parameter~3
(values $7/6/5$, $7/7/6$, $8/7/7$, and $r+1$ throughout for $r\ge7$).  The
initial target, Xie--Xu Conjecture~4.8 (arXiv:2408.09152v1), had already been
disproved by Alex Chengyu Li in September 2026; the report states this and
refines the target to the rankwise minimum-size question.}

\entry{A sharp order-five threshold for higher-order Stirling subset log-concavity}
{log-concavity-and-unimodality/higher-order-stirling-rows}
{article.pdf}{stirling\_log\_concavity.zip}
{For $S^{(r)}_{n,k}$, the number of partitions of $n+(r-1)k$ labels into $k$
blocks of size ${\ge}\,r$, proves that all rows are log-concave if and only if
$1\le r\le5$, with strict inequalities in the interior --- the log-concavity
assertion of Deb--Sokal Conjecture~1.4(c) (arXiv:2507.18959v1).  Also proves the
cycle cases $r=3,4$; the fifth-order \emph{cycle} conjecture is explicitly left
unresolved, and the included operator counterexample is flagged as not being a
counterexample to it.}

\entry{A sum-of-squares proof of Ulas's real-zero conjecture}
{log-concavity-and-unimodality/stern-polynomial-positivity}
{stern\_conjecture\_note.pdf}{stern\_conjecture\_research.zip}
{Conjecture~4.4 of Maciej Ulas, arXiv:1909.10844v1, for the Stern polynomials
indexed by $a_n=(4^n-1)/3$ and $h_n=2(4^n-1)(2\cdot4^n+1)/3+1$: for every
$n\ge0$ and every real $t$,
$W_n(t)=(A_{n+1}(t)-tA_n(t)/2)^2+3t^2A_n(t)^2/4>0$.  The adjacent-polynomial
quadratic identity behind it was already in Ulas's paper and is re-derived, not
claimed as new.}

\section{Graph theory}

\entry{Mutual-visibility sets in Andr\'asfai graphs}
{graph-theory/andrasfai-mutual-visibility}
{article.pdf}{andrasfai\_visibility\_research.zip}
{Proves the generating function attributed to Andrew Howroyd in OEIS A391632,
$\sum a(n)x^n = x(4-19x+33x^2-40x^3+8x^4)/(1-10x+29x^2-32x^3+8x^4)$, by an
exact weight-preserving bijection between mutual-visibility sets of the
$n$-Andr\'asfai graph and anchored closed walks of length $3n-1$ in a
four-state labelled graph.  Adds the full bivariate generating function, an
exact binomial formula for every cardinality, maximum-set classification,
eventual polynomiality at each fixed deficit, exponential asymptotics, a
central limit theorem, and linear-time recognition and exact uniform sampling.}

\entry{A bicyclic 30-vertex graph with domination root $-4$}
{graph-theory/domination-root-minus-four-order-30}
{domination\_root\_30.pdf}{domination\_root\_30\_research.zip}
{An explicit connected planar graph on 30 vertices and 31 edges, treewidth 2 and
domination number 10, with $D(G,x)=x^{10}(x+4)Q_{19}(x)$ --- a negative answer
to question~2 of Section~4 in Alikhani and Griswold, arXiv:2608.00109v1, which
asks whether 33 is the minimum order.  Their 33-vertex example had already
disproved the original integer-root conjecture; no claim of global minimality
at 30 is made.  Three independent exact polynomial computations agree.}

\entry{A marked-cycle proof of the trapezohedral domination conjecture}
{graph-theory/trapezohedral-minimal-dominating-sets}
{article.pdf}{trapezohedral\_domination\_research.zip}
{Proves the generating function labelled conjectural in OEIS A381190: for the
$n$-trapezohedral graph, the number of dominating sets that are
inclusion-minimal among all dominating sets and induce a connected subgraph is
$a_n=2n(2p_{n-2}+p_{n-3})$ with $p_m=p_{m-2}+p_{m-3}$ (Padovan).  The proof
classifies such sets by a cyclic arrangement of gaps of length 2 and 3 with a
restricted marked edge.  Adds the size distribution, strict log-concavity, an
unbiased sampler, a Binet formula and a central limit theorem.}

\section{Enumerative combinatorics}

\entry{A Twelve-Value Proof of the A181280 Formula}
{enumerative-combinatorics/a181280-binary-matrix-formula}
{article.pdf}{A181280\_research\_package.zip}
{The binary-matrix formula recorded as Conjecture~20 by Kauers and Koutschan
(2023) and still marked conjectural in OEIS A181280: an explicit spectral
recurrence bound and a finite certificate resting on twelve decisive counts,
together with further consequences.  \texttt{STATUS.md} records that general
C-finiteness was already proved in Kauers's December 2025 lecture slides.}

\entry{Esthetic Numbers, Folded Paths, and Algebraic Diagonals}
{enumerative-combinatorics/esthetic-numbers-and-folded-paths}
{article.pdf}{esthetic\_walks\_research.zip}
{The diagonal conjecture $T(q,q-2)=A182555(q-2)$ in OEIS A377000, where
$T(q,k)$ counts base-$q$ positive integers of length $k$ whose adjacent digits
differ by one.  Proves it exactly, proves the entry's other two conjectures as
retrieved, and derives formulas for every fixed-offset diagonal $q=k+d$.}

\entry{Diagonal Hardinian Arrays: A Fixed-Parameter Asymptotic Formula}
{enumerative-combinatorics/hardinian-array-diagonal-asymptotics}
{article.pdf}{hardinian\_arrays\_research.zip}
{For every fixed $r\ge0$, $H_r(n,n)\sim C_r4^{rn}n^{-r(r-1)/2}$ with an explicit
$C_r$ in Gamma values.  The formula confirms the constants tabulated by
Dougherty-Bliss and Kauers (2024) but refutes their proposed restriction to
powers of 2, 3 and $\pi$: $C_7=2^{18}5^2/(3^{45}\pi^3)$.  Adds an exact finite
determinant at fixed array size, an exact enumeration algorithm, and a limiting
negative-binomial law for total deficits.}

\entry{An explicit formula for lucky parking spots}
{enumerative-combinatorics/lucky-parking-spots}
{article.pdf}{lucky\_parking\_spots\_research.zip}
{Proves Conjecture~13 of Butler, Hadaway, Lenius, Martens and Moats (JIS 29
(2026), Article 26.1.1): a closed form for $C(n,j)$, the number of parking
functions in which \emph{spot} $j$ is lucky.  Adds an exact reflection law, all
fixed right-edge columns, a closed form for OEIS A374533, Poisson edge limits,
uniform bulk convergence, polynomial reciprocity and asymptotic expansions.}

\entry{Strict growth and the Catalan limit for adjacency-bounded 132-avoiding permutations}
{enumerative-combinatorics/adjacency-bounded-132-avoiders}
{article.pdf}{strict\_growth\_132.zip}
{For $132$-avoiding permutations with every adjacent difference at most $m$ and
growth constant $\alpha_m$: proves $\alpha_{m+1}>\alpha_m$ for all $m\ge1$ ---
the remaining strict-inequality part of Nadler's Conjecture~2 --- and the
refinement $4-\alpha_m=(2\log m+4\log\log m+O(1))/m$, with
$2\log\pi-4\log2\le\liminf E_m\le\limsup E_m\le2\log\pi$.  The all-$m$
component ordering and convergence of $E_m$ are explicitly not claimed.}

\entry{Throwback Dynamics with Repeated Entries}
{enumerative-combinatorics/throwback-queue-dynamics}
{throwback\_research.pdf}{throwback\_research.zip}
{For the queue process that removes the head and reinserts it at the zero-based
position equal to its weight: every token recurs if and only if at most $m$
input tokens have weight at most $m$, for every $m\ge1$, with an explicit
description of the recurrent tokens when that fails.  Adds exact cycle periods
and frequencies, a coprimality criterion for perfect mixing, Catalan and
prime-parking counts for recurrent cores, and the largest labelled period of a
core of size $r$.}

\entry{Reciprocity and Matrix Duality for Preorder Polytopes}
{enumerative-combinatorics/preorder-polytope-reciprocity}
{article.pdf}{preorder\_reciprocity.zip}
{Proves Conjectures 7.2, 4.2 and 4.4 of Athanasiadis and Chapoton,
\emph{Polytopes and posets associated to preorders} (arXiv:2605.26916v1), and
the ordinary $q=1$ case of their Conjecture~4.9 (not its full $q$-refinement).
Two mechanisms: Ehrhart--Macdonald reciprocity after subtracting the all-ones
vector, extending the double identity to independently weighted capacities; and
duplication of both colour classes of the bipartite graph followed by
Postnikov's lattice-point formula, giving the conjectured matrix transpose.}

\entry{Rationality and growth of stabilized permutation-weight series}
{enumerative-combinatorics/stabilized-permutation-weights}
{article.pdf}{stabilized\_permutation\_weights.zip}
{Studies the full stabilized generating functions of the permutation-weight
$q$-Eulerian polynomials, beyond their initial connection with OEIS A256193:
an all-$d$ finite-spectrum theorem, rationality, a universal integer
denominator, the exact radius of convergence and residue-class asymptotics.
The two-descent sequence gets an explicit formula in powers of 3, powers of 2
and Fibonacci numbers.}

\section{Generating functions and asymptotics}

\subsection{Coefficient counts of binary products}

Four reports on counting coefficients of $\prod_i(1+x^{G_i})$.  The first three
answer, independently, the universal-rationality part of Richard Stanley's
MathOverflow question 431075 (``A conjectured rational generating function'',
2022), which had no posted answer when checked on 19 September 2026.

\entry{When Rational Exponents Produce Non-D-Finite Coefficient Counts}
{generating-functions-and-asymptotics/binary-product-coefficient-counts/mixed-base-natural-boundary}
{article.pdf}{Stanley\_Rationality\_Counterexample.zip}
{With $G(2j+1)=2^j$ and $G(2j+2)=3^j$, the exponent generating function is
rational but the generating function of $\nu(N)$ --- the number of coefficients
equal to one in $P_N(x)=\prod_{i\le N}(1+x^{G(i)})$ --- is not D-finite and has
a natural boundary; explicitly,
$a(n)=\nu(2n)=2^{\lfloor(n+1)(1-\log2/\log3)\rfloor}$ for $n\ge2$.  The proof
assumes nothing about the algebraicity of the growth constant.  In the family
$G_b$ the count series is rational exactly when $b$ is a power of two.}

\entry{A Mixed-Base Counterexample to Stanley's Rationality Question}
{generating-functions-and-asymptotics/binary-product-coefficient-counts/mixed-base-rationality-classification}
{article.pdf}{stanley\_mixed\_base\_counterexample.zip}
{The same mixed-base construction, with
$\nu(2n)=2^{\,n-\lfloor(n+1)\alpha\rfloor}$ for $\alpha=\log2/\log3$ proved by
interval coverage and a carry-gap count.  Nonrationality is obtained by a
finite-field argument, then strengthened to a natural boundary,
non-D-finiteness and non-P-recursiveness.  Classifies rationality in the family
interleaving $a^j$ and $q^j$ with digit polynomial $1+y+\cdots+y^{a-1}$: it
holds exactly when $q^u=a^v$.}

\entry{Rational exponents, nonrational coefficient counts}
{generating-functions-and-asymptotics/binary-product-coefficient-counts/recursive-exponent-construction}
{research.pdf}{stanley\_counterexample.zip}
{A different negative answer, from the recursion $a_0=2$, $a_1=1$,
$a_n=a_{n-1}-2a_{n-2}$, with $b_n=T_{n-1}+2^n-1+a_n$, $T_n=T_{n-1}+b_n$,
$G_{2n-1}=2^{n-1}$ and $G_{2n}=b_n$: the counting series has a natural boundary
on the unit circle, hence is neither rational, algebraic nor D-finite.}

\entry{Odd coefficients in quasifibonacci products}
{generating-functions-and-asymptotics/binary-product-coefficient-counts/quasifibonacci-odd-coefficient-counts}
{article.pdf}{quasifibonacci\_parity\_research.zip}
{Stanley's Conjecture~6.3, with a seed-independent extension.  For $k\ge2$ and
seeds each exceeding the sum of their predecessors, continued by the $k$-term
recurrence, the number $h_k(n)$ of odd coefficients of
$\prod_{i\le n}(1+x^{w_i})$ has generating function
$(1+2z^k)/(1-2z+2z^k-2z^{k+1})$, independently of the seeds.  The proof builds
periodic signs of block length $k+1$ and block product $-1$ and shows every
signed finite product is flat.}

\subsection{Asymptotics, records and discrepancy}

\entry{The Exponential Scale of Gregory Sign Thresholds}
{generating-functions-and-asymptotics/gregory-coefficient-sign-thresholds}
{gregory\_sign\_threshold.pdf}{gregory\_sign\_threshold\_research.zip}
{For $(x/\log(1+x))^m=\sum_n G^{(m)}_nx^n$ and $\mathrm{Gr}(m)$ the least $N$
beyond which $(-1)^{n-1}G^{(m)}_n>0$, proves
$\log\mathrm{Gr}(m)=m+1-\gamma-(\zeta(2)+\zeta(3))/m+C_2/m^2+O(m^{-3})$ with
$C_2$ explicit, plus a finite algorithm for every coefficient of the complete
expansion; in particular $\mathrm{Gr}(m)/e^m\to e^{1-\gamma}$ and
$\mathrm{Gr}(m)/3^m\to0$.  This replaces Conjecture~6.1 of Xu and Zhao
(arXiv:2609.11072v1), whose printed lower bound already conflicts with its own
reported $\mathrm{Gr}(5)=113$; that value is independently certified here.}

\entry{Exact Records in the Square-Root-of-Three Walk}
{generating-functions-and-asymptotics/sqrt-three-walk-records}
{article.pdf}{sqrt3\_record\_recurrence.zip}
{Proves the four-phase record recurrence for
$S(N)=\sum_{n\le N}(-1)^{\lfloor n\sqrt3\rfloor}$ in the corrected form on
page~7 of Bruin and Fokkink, \emph{On the records and zeros of a deterministic
random walk} (arXiv:2503.11734v2), where it is experimental and attributed to
Arias de Reyna and van de Lune.  Includes minimality of every claimed first
passage, a rational generating function, closed forms, sharp logarithmic
extrema, and an extension to every even integer trace $D\ge4$.}

\entry{Sharp discrepancy for a family of binary substitutions}
{generating-functions-and-asymptotics/binary-substitution-discrepancy}
{article.pdf}{oeis\_discrepancy\_counterexamples.zip}
{For the fixed word of $0\mapsto1$, $1\mapsto101010$, the position sequences
A284365 and A284366 both carry a conjectured upper position error below~2, and
both conjectures are false: the first counterexamples are at ranks
$2{,}977{,}771$ and $8{,}933{,}313$, with common error
$(2977771\sqrt{21}-13645857)/2\approx2.00487$ and sharp common upper bound
$(6+\sqrt{21})/5\approx2.11652$.  Proves complete discrepancy limit intervals
for every substitution $0\mapsto1$, $1\mapsto(10)^m$, with infinite
counterexample families and exact recurrences for their indices.}

\end{document}
